\documentclass[10pt,a4paper]{amsart}
\usepackage[margin=19mm]{geometry}
\usepackage{amsmath,amssymb,mathtools}
\usepackage[scr=rsfs,cal=euler]{mathalfa}
\usepackage{xfrac}
\usepackage[unicode,hidelinks,bookmarks=false]{hyperref}
\newcommand{\ie}{i.e.\ }
\newcommand{\cf}{cf.\ }
\newcommand{\resp}{respectively\ }
\newcommand{\Subsection}{\subsection}
\providecommand{\nobreakdash}{\nobreak\hbox{-}}
\setlength{\emergencystretch}{2em}
\multlinegap0pt
\usepackage{xcolor}
\usepackage{amssymb}
\usepackage{mathtools}
\usepackage{dsfont}
\usepackage{cancel}
\usepackage[shortlabels]{enumitem}
\usepackage{float}
\usepackage{bm}
\usepackage{tikz}
\usepackage{algorithm}\usepackage{algpseudocode}
\newtheorem{theorem}{Theorem}
\newtheorem{prop}{Proposition}
\newtheorem{lemma}{Lemma}
\newcommand{\Crad}{\cC_{\rm rad}}
\newcommand*{\E}{\mathbb E}
\DeclareMathOperator*{\KL}{KL}
\newcommand{\R}{\mathbb{R}}
\DeclareMathOperator*{\argmin}{arg\,min}
\newcommand{\cC}{\mathcal{C}}
\newcommand{\cF}{\mathcal{F}}
\newcommand{\cT}{\mathcal{T}}
\newcommand*{\dd}{\, \rd}
\newcommand*{\defeq}{\coloneqq}
\newcommand{\eps}{\varepsilon}
\renewcommand{\epsilon}{\varepsilon}
\newcommand{\kl}[2]{\KL(#1 \!\;\|\; \!#2)}
\newcommand{\pirad}{\pi_{\rm rad}}
\newcommand*{\rd}{\mathrm{d}}
\renewcommand{\tilde}{\widetilde}
\DeclareMathOperator{\tr}{tr}
\title{Variational inference via radial transport}
\author{Luca Ghafourpour, Sinho Chewi, Alessio Figalli, Aram-Alexandre Pooladian}
\date{Source: arXiv:2602.17525v2 (CC BY 4.0)}
\begin{document}
\maketitle

\noindent\emph{Source and licence.} Variational inference via radial transport. Version arXiv:2602.17525v2 (CC BY 4.0), distributed by arXiv under the Creative Commons Attribution 4.0 International licence, http://creativecommons.org/licenses/by/4.0/, as recorded in the arXiv licence metadata for this exact version and shown on the abstract page https://arxiv.org/abs/2602.17525v2. Full licence notice: \texttt{licenses/arxiv-2602.17525-CC-BY-4.0.txt}.\par\smallskip

\noindent\textit{Editorial scope.} Counted unit: the article's lemma lem:sgd\_lemma2 (source0.tex lines 1244--1250). The article's complete original proof of it is delivered with it (source0.tex lines 1251--1269, one \texttt{proof} environment of 19 lines). No mathematics is written, repaired or re-derived: the statement and the proof are the article's own lines, in the article's own notation, and no retained line is substituted or re-flowed.  Retained uncounted as the source-local context the proof needs: lines 301--303; lines 307--309; lines 345--354; lines 387--389; lines 396--402; lines 491--493.  Nothing was omitted from any retained span, so the package carries no omission record.  No retained line contains a citation, so the wrapper carries no bibliography and no bibliography entry.  No retained line contains a figure environment, an image-inclusion command or drawing code, so no image is delivered and the package declares no asset dependency.  Editorial formatting only, in addition: an AMS \texttt{amsart} preamble that restates the article's own declarations and macros used by the retained lines, namely the environments \texttt{theorem}, \texttt{prop}, \texttt{lemma} on separate counters, together with the standard packages those lines need.  The retained environments print in this document as Theorem 1, Proposition 1, Lemma 1 in order of appearance, whereas the article prints the same items with the numbering of its own full document.  The complete native source of the article as distributed in the arXiv e-print for this exact version is bundled unchanged as \texttt{sources/194969/source0.tex}.
\begin{align}\label{eq:radvi}\tag{$\mathsf{radVI}$}
    \pi_{\rm rad} \in \argmin_{\mu \in \Crad} \kl{\mu}{\pi}\,,
\end{align}

\begin{align}\label{well_cond} \tag{$\mathsf{WC}$}
0 \preceq \ell_V I \preceq \nabla^2 V \preceq L_V I\,,
\end{align} 

\begin{theorem}[Regularity of the optimal radial map]\label{theorem_map_regularity} 
Suppose $\pi$ satisfies \eqref{well_cond} and consider the corresponding minimizer to \eqref{eq:radvi}, denoted $\pirad^\star$. Let $T^\star_{\rm{rad}}$ denote the unique optimal transport map from $\rho$ to $\pirad^\star$, and write $T_{\rm rad}^\star(x) = \Psi^\star(\|x\|)\,x/\|x\|$. Writing $r = \|x\|$, the following regularity conditions hold:
\begin{align*}
    \frac{1}{\sqrt{L_V}} \leq (\Psi^\star)'(r) &\leq \frac{1}{\sqrt{\ell_V}}\,, \\
    \left|(\Psi^\star)''(r) \right|
    &\lesssim
    {\frac{\kappa}{\sqrt{\ell_V}}\,\bigl(1 + \frac{d}{r^2}\bigr)\,(1+|r-\sqrt d|)}\,,
\end{align*}
where $\kappa \defeq L_V/\ell_V$.
\end{theorem}

\begin{align}\label{eq:basis_functions}
    {\Psi}_j(r) \defeq {\Psi}_{\rm base}(\delta_j^{-1}(r-a_j))\,,
\end{align}

\begin{theorem}[Universal approximation]\label{theorem_unif_approx}
Let $T^\star \in \cT_{\rm rad}$ which satisfies Theorem~\ref{theorem_map_regularity}, and let ${\epsilon \gg \exp(-\Omega(d))}$. Define $\cT_J$ with $\alpha = 1/\sqrt{L_V}$ and choose $R \asymp {\sqrt{\log(d/\epsilon)}}$, $J = {\widetilde{\Omega}(\sqrt{\kappa/\eps})}$ with the basis elements given by \eqref{eq:basis_functions}. Then there exists a $\hat{T} \in \cT_J$, i.e., there exists $\hat{\lambda} \in \R^{J+1}_+$ with $\hat{T} = T_{\hat{\lambda}}$, such that
\begin{align*}
    \|T^{\star} - \hat{T}\|_{L^2(\rho)} &\leq \epsilon/\ell_V^{1/2}\,,\\
    \|D(T^{\star} - \hat{T} ) \|_{L^2(\rho)} &\le {\widetilde{\mathcal{O}}(\kappa^{1/2} \epsilon^{1/2}/\ell_{V}^{1/2})}\,.
\end{align*}
\end{theorem}

\begin{prop}\label{prop:cf_smooth_convex}
Suppose $\pi$ satisfies \eqref{well_cond}, and also consider $\cT_J$ chosen as in Theorem~\ref{theorem_unif_approx}. Then $\lambda\mapsto \cF(T_\lambda)$ is $\ell_V$-strongly convex and ${\widetilde O(J^2 L_V)}$-smooth (with respect to $\|\cdot\|_Q$).
\end{prop}

\begin{lemma}\label{lem:sgd_lemma2}
Let $\{\Psi_j\}_{j=0}^J$ be our dictionary of choice, and let $Q_{ij} = \E_{X\sim\rho}[\Psi_i(\|X\|)\Psi_j(\|X\|)]$. Then for any $x \in \R^d$, it holds that
\begin{align*}
    \langle Q^{-1}, \vec\Psi(x)\vec\Psi(x)^\top\rangle \lesssim M^3.
\end{align*}
where $M = J+1$.
\end{lemma}

\begin{proof}
Our goal is to show that ${Q}$ is lower bounded by some absolute constant $c_Q > 0$. This is sufficient, as
\begin{align*}
    \langle Q^{-1}, \vec\Psi(x)\vec\Psi(x)^\top \rangle  {\le} \frac{1}{c_Q}\tr(\vec\Psi(x)\vec\Psi(x)^\top) = \frac{1}{c_Q}\sum_{j=0}^J\Psi_j^2(\|x\|) \leq \frac{M}{c_Q}\,,
\end{align*}
where we used the trivial upper bound $\Psi_j \leq 1$. Ultimately, we want to show that $1/c_Q \lesssim M^2$. Taking $\lambda \in \R^M_+$, this is equivalent to proving the following lower bound
\begin{align*}
    \lambda^\top Q \lambda = \int \Bigl(\sum_{j=0}^J \lambda_j\Psi_j(r)\Bigr)^2 \dd \tilde\rho(r) \eqqcolon \int (\tilde T_\lambda(r))^2 \dd \tilde\rho(r) \overset{!}{\gtrsim} M^{-2}\|\lambda\|^2\,.
\end{align*}
We can lower-bound this quantity term by term along the partitions $I_0 = [0,\sqrt{d}-R]$ and $I_\ell = [a_\ell,a_\ell + \delta]$. For $\ell \geq 1$,
\begin{align*}
    \int_{a_\ell}^{a_\ell + \delta} (\tilde T_\lambda(r))^2 \dd \tilde\rho(r) \geq \lambda_\ell^2 \inf_m \int_{a_\ell}^{a_\ell + \delta} ((r - a_\ell)/\delta - m)^2 \dd \tilde\rho(r) \gtrsim \lambda_\ell^2 \Upsilon^{-1}\int_{a_\ell}^{a_\ell + \delta}\dd\tilde\rho(r)\gtrsim \lambda_\ell^2 \Upsilon^{-1}\,,
\end{align*}
where, over the interval, $\tilde\rho$ is effectively constant (recall the arguments in Proposition~\ref{prop:cf_smooth_convex}). Also, for $\ell=0$, we can use the arguments from Proposition~\ref{prop:cf_smooth_convex}) in this special case to prove that
\begin{align*}
    \int_{0}^{\delta_0} (\tilde T_\lambda(r))^2 \dd \tilde\rho(r) \geq \lambda_0^2 \delta_0^{-2} \int_0^{\delta_0} r^2 \dd \tilde\rho(r) \gtrsim \lambda_0^2 \,.
\end{align*}
Adding up all the terms, this concludes the proof as $\Upsilon \asymp M^2$.
\end{proof}

\end{document}
